% !PRINTED 442
\weberpage{442}
% running head: 442  Achtzehnter Abschnitt.  § 122.

%<<CONT
\gesperrt{und da $\mathfrak{S}$ eine Abelsche Gruppe ist, so ist die Klassengleichung eine Abelsche.}

Nehmen wir an, die Funktion $R$ habe reelle (rationale) Koeffizienten, so wird gleichwohl ihr rationaler Ausdruck die Form $a+b\sqrt{D}$ haben, worin $a, b$ rationale Zahlen sind. Der imaginäre Teil wird dann und nur dann wegfallen, wenn $R$ auch durch die Vertauschung sämtlicher Klassen $k, k', k'', \ldots$ mit ihren entgegengesetzten ungeändert bleibt. Daraus folgt, daß ohne Adjunktion von $\sqrt{D}$ die \gesperrt{Galoissche} Gruppe der Klassengleichung enthalten ist in dem System von $2h$ Permutationen, die man erhält, wenn man in $\mathfrak{S}$ jede Klasse in ihre entgegengesetzte verwandelt. Nur in dem besonderen Falle, in dem alle Klassen zweiseitig sind (der nur für eine endliche Anzahl von Determinanten stattfindet), ist dies letztere System mit $\mathfrak{S}$ identisch, und die Klassengleichung ist ohne Adjunktion von $\sqrt{D}$ eine \gesperrt{Abelsche}.

Der algebraische Körper, der aus den rationalen Funktionen einer Klasseninvariante gebildet ist, ist daher, von dem zuletzt erwähnten Ausnahmefall abgesehen, kein \gesperrt{Normalkörper}, sondern ist von seinen konjugierten Körpern verschieden. Dagegen erhält man einen Normalkörper, wenn man die Quadratwurzel $\sqrt{D}$ dem Körper der Klasseninvarianten adjungiert. Denn jede rationale Funktion sämtlicher Wurzeln der Gleichung kann durch eine von ihnen und $\sqrt{D}$ rational ausgedrückt werden.

\begin{center}
§ 122. \gesperrt{Irreducibilität.}
\end{center}

Wir betrachten jetzt den algebraischen Zahlkörper
\weq{1}{\mathfrak{K} = \mathfrak{R}(k, \sqrt{D}),}
der aus einer Klasseninvariante $(k)$ der Diskriminante $D$ und $\sqrt{D}$ zusammengesetzt ist; in diesem Körper sind nach dem vorigen Paragraphen alle zu derselben Diskriminante gehörenden Klasseninvarianten $(k), (k'), (k''), \ldots$ enthalten. Diesen Körper nennen wir den \gesperrt{Klassenkörper} der Diskriminante $D$. Den quadratischen Körper $\mathfrak{R}(\sqrt{D})$ bezeichnen wir wie bisher mit $\Omega$, der, wenn $D=Q^2\Delta$ und $\Delta$ der Stamm von $D$ ist, mit $\mathfrak{R}(\sqrt{\Delta})$ identisch ist.
$h$ sei die \gesperrt{Klassenzahl} der Diskriminante $D$.

1. Es sei $(k)$ irgend eine der $h$ Klassenvarianten und % [sic] printed "Klassenvarianten"; context requires "Klasseninvarianten"
\weq{2}{\Phi(k) = (k)^\nu + \alpha_1(k)^{\nu-1} + \alpha_2(k)^{\nu-2}\cdots + \alpha_\nu = 0}
die Gleichung niedrigsten Grades in $\Omega$, deren Wurzel $(k)$ ist. Da $(k)$ eine ganze Zahl ist, so müssen auch die $\alpha_i$ ganze Zahlen
%CONT>>
